% config/fonts.tex — 字体配置（XeLaTeX + ctex fontset=none）
% ==========================================================

% IfFontExistsTF fallback（兼容旧版 fontspec）
\ifdefined\IfFontExistsTF\else
  \newcommand{\IfFontExistsTF}[3]{#3}
\fi

% \xeCJKsetup{PunctStyle=plain}

% ---------- 英文正文（Palatino 风格，使用 TeX Live OTF）----------
\setmainfont{texgyrepagella}[
  Extension      = .otf,
  UprightFont    = *-regular,
  BoldFont       = *-bold,
  ItalicFont     = *-italic,
  BoldItalicFont = *-bolditalic,
]
\setsansfont{texgyreheros}[
  Extension      = .otf,
  UprightFont    = *-regular,
  BoldFont       = *-bold,
  ItalicFont     = *-italic,
  BoldItalicFont = *-bolditalic,
]
\setmonofont{Menlo}

% ---------- 数学字体 ----------
\setmathfont{texgyrepagella-math.otf}

% ---------- 中文正文 —— 思源字体 ----------
\IfFileExists{\AMaNRoot _asset/fonts/SourceHanSerifSC-Regular.otf}{%
  \setCJKmainfont{SourceHanSerifSC-Regular.otf}[
    Path           = {\AMaNRoot _asset/fonts/},
    BoldFont       = SourceHanSerifSC-Bold.otf,
    ItalicFont     = SourceHanSerifSC-Regular.otf,
    BoldItalicFont = SourceHanSerifSC-Bold.otf,
  ]
}{%
  \setCJKmainfont{Source Han Serif SC}[
    BoldFont       = Source Han Serif SC Bold,
    ItalicFont     = Source Han Serif SC,
    BoldItalicFont = Source Han Serif SC Bold,
  ]
}
\IfFileExists{\AMaNRoot _asset/fonts/SourceHanSansSC-Regular.otf}{%
  \setCJKsansfont{SourceHanSansSC-Regular.otf}[
    Path     = {\AMaNRoot _asset/fonts/},
    BoldFont = SourceHanSansSC-Bold.otf,
  ]
  \setCJKmonofont{SourceHanSansSC-Regular.otf}[
    Path = {\AMaNRoot _asset/fonts/},
  ]
}{%
  \setCJKsansfont{Source Han Sans SC}[
    BoldFont = Source Han Sans SC Bold,
  ]
  \setCJKmonofont{Source Han Sans SC}
}

% ---------- 中文黑体 ----------
\IfFileExists{\AMaNRoot _asset/fonts/SourceHanSansSC-Bold.otf}{%
  \setCJKfamilyfont{zhhei}{SourceHanSansSC-Bold.otf}[
    Path = {\AMaNRoot _asset/fonts/},
  ]
}{%
  \setCJKfamilyfont{zhhei}{Source Han Sans SC Bold}
}
\newcommand{\heiti}{\CJKfamily{zhhei}}

% ---------- 中文楷体（最小定义，供 \kaiti 使用） ----------
\setCJKfamilyfont{zhkai}{FandolKai-Regular.otf}
\newcommand{\kaiti}{\CJKfamily{zhkai}}

% ---------- 等宽文字（与 URL 共用字号） ----------
\newcommand{\TTTFont}{\fontsize{9.5pt}{\baselineskip}\selectfont\ttfamily}
\DeclareRobustCommand{\ttt}[1]{{\TTTFont #1}}

% ==========================================
% 全角中文标点统一压缩
% 仍然直接输入：（），《》，“”，，。！？：；、……
% 编译时自动应用新的左右间距
% 需要：XeLaTeX 或 LuaLaTeX
% ==========================================

\usepackage{newunicodechar}

% ------------------------------------------
% 统一可调参数
% ------------------------------------------
\newdimen\CJKPunctOuterKern   % 标点外侧压缩
\newdimen\CJKPunctInnerKern   % 成对标点内侧压缩
\newdimen\CJKDashKern         % 破折号两侧统一间距

% 你主要调这两个值
\CJKPunctOuterKern=-0.2em
\CJKPunctInnerKern=+0.08em
\CJKDashKern=-0.06em

% ------------------------------------------
% 工具宏
% ------------------------------------------
\makeatletter

% 单体标点：左右都用 outer
\newcommand{\CJKPunctSingle}[1]{%
  \nobreak\hspace*{\CJKPunctInnerKern}#1\hspace*{\CJKPunctOuterKern}%
}

% 破折号单独处理，避免左右间距不对称
\newcommand{\CJKDash}[1]{%
  \nobreak\hspace*{\CJKDashKern}#1\hspace*{\CJKDashKern}%
}

% 左半边成对标点：外侧 outer，内侧 inner
\newcommand{\CJKPunctOpen}[1]{%
  \nobreak\hspace*{\CJKPunctOuterKern}#1\nobreak\hspace*{\CJKPunctInnerKern}%
}

% 右半边成对标点：内侧 inner，外侧 outer
\newcommand{\CJKPunctClose}[1]{%
  \nobreak\hspace*{\CJKPunctInnerKern}#1\hspace*{\CJKPunctOuterKern}%
}

% 中文省略号需要把“……”当成一个整体处理，不能把两个“…”分别加间距。
\newcommand{\CJKEllipsisSingle}{%
  \nobreak\hspace*{\CJKPunctInnerKern}\char"2026\hspace*{\CJKPunctOuterKern}%
}
\newcommand{\CJKEllipsisPairOutput}{%
  \nobreak\hspace*{\CJKPunctInnerKern}\char"2026\char"2026\hspace*{\CJKPunctOuterKern}%
}
\begingroup
\lccode`\~=`\…
\lowercase{%
\endgroup
\def\CJKEllipsis{%
  \futurelet\nexttoken\CJKEllipsisAux
}
\def\CJKEllipsisAux{%
  \ifx\nexttoken~
    \expandafter\CJKEllipsisPair
  \else
    \CJKEllipsisSingle
  \fi
}
\def\CJKEllipsisPair~{%
  \CJKEllipsisPairOutput
}
}

\makeatother

% ------------------------------------------
% 成对标点：圆括号 / 方头括号 / 方括号 / 花括号
% ------------------------------------------
\newunicodechar{（}{\CJKPunctOpen{\char"FF08}}
\newunicodechar{）}{\CJKPunctClose{\char"FF09}}

\newunicodechar{【}{\CJKPunctOpen{\char"3010}}
\newunicodechar{】}{\CJKPunctClose{\char"3011}}

\newunicodechar{〔}{\CJKPunctOpen{\char"3014}}
\newunicodechar{〕}{\CJKPunctClose{\char"3015}}

\newunicodechar{［}{\CJKPunctOpen{\char"FF3B}}
\newunicodechar{］}{\CJKPunctClose{\char"FF3D}}

\newunicodechar{｛}{\CJKPunctOpen{\char"FF5B}}
\newunicodechar{｝}{\CJKPunctClose{\char"FF5D}}

% ------------------------------------------
% 成对标点：书名号 / 尖括号
% ------------------------------------------
\newunicodechar{《}{\CJKPunctOpen{\char"300A}}
\newunicodechar{》}{\CJKPunctClose{\char"300B}}

\newunicodechar{〈}{\CJKPunctOpen{\char"3008}}
\newunicodechar{〉}{\CJKPunctClose{\char"3009}}

% ------------------------------------------
% 成对标点：中文引号
% ------------------------------------------
\newunicodechar{「}{\CJKPunctOpen{\char"300C}}
\newunicodechar{」}{\CJKPunctClose{\char"300D}}

\newunicodechar{『}{\CJKPunctOpen{\char"300E}}
\newunicodechar{』}{\CJKPunctClose{\char"300F}}

\newunicodechar{“}{\CJKPunctOpen{\char"201C}}
\newunicodechar{”}{\CJKPunctClose{\char"201D}}

\newunicodechar{‘}{\CJKPunctOpen{\char"2018}}
\newunicodechar{’}{\CJKPunctClose{\char"2019}}

% ------------------------------------------
% 单体标点：顿号、逗号、句号、分号、冒号、问号、叹号、省略号
% ------------------------------------------
\newunicodechar{、}{\CJKPunctSingle{\char"3001}}
\newunicodechar{，}{\CJKPunctSingle{\char"FF0C}}
\newunicodechar{。}{\CJKPunctSingle{\char"3002}}
\newunicodechar{；}{\CJKPunctSingle{\char"FF1B}}
\newunicodechar{：}{\CJKPunctSingle{\char"FF1A}}
\newunicodechar{？}{\CJKPunctSingle{\char"FF1F}}
\newunicodechar{！}{\CJKPunctSingle{\char"FF01}}
\newunicodechar{…}{\CJKEllipsis}

% ------------------------------------------
% 可选：连接号 / 破折号
% 这一类视觉差异较大，先给统一版本
% ------------------------------------------
\newunicodechar{—}{\CJKDash{\char"2014}}
\newunicodechar{～}{\CJKPunctSingle{\char"FF5E}}
